\documentclass[10pt]{article}
\usepackage[a4paper,margin=25mm]{geometry}
\usepackage{amsmath,amssymb,amsthm}
\usepackage[hidelinks]{hyperref}
\emergencystretch=2em
\newtheorem{theorem}{Theorem}
\title{The exact discrete gradient of the average-vector-field method\\
\large Proof of Conjecture 00000007367}
\author{GitHub: gaochengzhecpu}\date{}
\begin{document}\maketitle
\begin{abstract}
The interval average of the differential of a continuously differentiable
energy is an exact discrete differential. It satisfies both the energy
difference identity and consistency on the diagonal. Pairing this average
with a skew Hamiltonian operator proves exact energy preservation of
every average-vector-field step. We prove the full identity for arbitrary
real normed spaces using genuine Fr\'echet derivatives and Bochner
integrals, and verify it in Lean.
\end{abstract}

\section*{Precise statement and conventions}
The bilingual source asserts the exact discrete-gradient identity and
its realization by the averaging operator in the average-vector-field
(AVF) method. We use its standard differentiable setting: an energy
$H$ with continuous first derivative, and exact integration in the
definition of the method. No restriction to quadratic energies is needed.

Let $E$ be a real normed vector space and let $E^*$ denote its continuous
dual. Suppose $H:E\to\mathbb R$ is differentiable and
$DH:E\to E^*$ is continuous. For $x,y\in E$, define
\[
 \gamma_{x,y}(t)=x+t(y-x),\qquad
 \overline{D}H(x,y)=\int_0^1 DH(\gamma_{x,y}(t))\,dt\in E^*.
\]
The integral is the Bochner integral. The dual $E^*$ is complete even
when $E$ is not, and the integrand is continuous on the compact interval,
so this integral exists.

In a finite-dimensional Euclidean space the Riesz identification gives
the familiar averaged gradient
\[
 \overline\nabla H(x,y)=\int_0^1\nabla H(x+t(y-x))\,dt,
 \qquad
 \overline{D}H(x,y)[v]=\langle\overline\nabla H(x,y),v\rangle.
\]
Thus the covector formulation below is the ordinary AVF discrete
gradient expressed without choosing coordinates or an inner product.

\begin{theorem}[Exact discrete differential]
For all $x,y\in E$,
\begin{align}
 \overline{D}H(x,y)[y-x]&=H(y)-H(x),\label{identity}\\
 \overline{D}H(x,x)&=DH(x).\label{consistency}
\end{align}
\end{theorem}
\begin{proof}
The chain rule gives, for every real $t$,
\[
 \frac{d}{dt}H(\gamma_{x,y}(t))
   =DH(\gamma_{x,y}(t))[y-x].
\]
This derivative is continuous. Evaluation at $y-x$ is a continuous linear
functional on $E^*$, hence commutes with the Bochner integral. The
fundamental theorem of calculus therefore yields
\begin{align*}
 \overline{D}H(x,y)[y-x]
 &=\int_0^1DH(\gamma_{x,y}(t))[y-x]\,dt\\
 &=H(\gamma_{x,y}(1))-H(\gamma_{x,y}(0))
  =H(y)-H(x).
\end{align*}
When $y=x$, the path is constant and the integral over an interval of
length one is $DH(x)$, proving consistency.
\end{proof}

\begin{theorem}[Exact AVF energy preservation]
Let $J:E^*\to E$ satisfy
\begin{equation}\label{skew}
 g[J(g)]=0\qquad\hbox{for every }g\in E^*.
\end{equation}
If $h\in\mathbb R$ and $x,y\in E$ satisfy the AVF step
\begin{equation}\label{step}
 y-x=hJ(\overline{D}H(x,y)),
\end{equation}
then $H(y)=H(x)$.
\end{theorem}
\begin{proof}
Put $g=\overline{D}H(x,y)$. Equations \eqref{identity}, \eqref{step},
and \eqref{skew} give
\[
 H(y)-H(x)=g[y-x]=g[hJ(g)]=h\,g[J(g)]=0.
\]
\end{proof}

For the usual Hamiltonian system $\dot z=S\nabla H(z)$ in
$\mathbb R^d$, take a constant skew-symmetric matrix $S$ and identify
covectors with vectors. Then $J(g)=Sg$ satisfies \eqref{skew}, because
$g^T Sg=-g^T Sg$. Linearity of $S$ makes \eqref{step} precisely
\[
 y-x=h\int_0^1 S\nabla H(x+t(y-x))\,dt,
\]
the average-vector-field method. Its energy conservation follows for
every energy under the stated differentiability hypotheses and every
step that solves the implicit equation. The theorem asserts preservation,
not existence or uniqueness of that implicit solution. Replacing the
integral by an inexact quadrature or solving the step inexactly requires
a separate error analysis and is not the exact method in the source.

\section*{Formal verification and scope}
The Lean development uses Mathlib's actual continuous linear maps,
\texttt{HasFDerivAt}, \texttt{fderiv}, and Bochner interval integral.
The main proposition universally quantifies over energies on any real
normed space, assumes differentiability and continuity of the actual
Fr\'echet derivative, and states the conjunction of the discrete-gradient
identity, diagonal consistency, and preservation of every AVF step with
\eqref{skew}. The main theorem is
\begin{center}\texttt{Conjecture7367.conjecture7367}.\end{center}
The proof establishes the chain rule along the actual affine segment,
commutes evaluation with the actual integral, and applies Mathlib's
fundamental theorem of calculus. No finite examples, numerical integration,
or surrogate energy-difference formula replace these objects.

The finite-dimensional Riesz identification and the short matrix argument
that skew-symmetry implies \eqref{skew} are presented on paper. Lean
proves the stronger coordinate-free preservation theorem directly under
\eqref{skew}. These are the only bridges from the formal theorem to the
usual gradient-and-skew-matrix notation. Both the identity and its
averaging realization are fully formalized, as is the entire conservation
argument. The allowed axioms printed by Lean are only
\texttt{propext}, \texttt{Classical.choice}, and \texttt{Quot.sound}.

\section*{Source and reproduction}
The unmodified bilingual statement is supplied as \texttt{SOURCE.md};
its frozen provenance and SHA-256 are recorded in the verification folder.
The coordinating workflow found no prior submission for this conjecture.
Lean is pinned to 4.19.0 and Mathlib commit
\begin{center}\texttt{c44e0c8ee63ca166450922a373c7409c5d26b00b}.\end{center}
In the supplied Lean folder run
\begin{center}\texttt{lake env lean -DwarningAsError=true Main.lean}.\end{center}
This draft was developed by a delegated AI agent with a self-review;
the coordinating agent reviews it separately. No external independent
review is claimed.
\end{document}
